\documentclass[11pt,a4paper]{report}
\usepackage[utf8]{inputenc}
\usepackage[T1]{fontenc}
\usepackage[french]{babel}
\usepackage{lmodern}
\usepackage{geometry}
\geometry{margin=2.5cm}
\usepackage{graphicx}
\usepackage{booktabs}
\usepackage{array}
\usepackage{amsmath}
\usepackage{siunitx}
\usepackage{hyperref}
\usepackage{float}
\usepackage{caption}
\usepackage{subcaption}
\usepackage{csquotes}
\usepackage{fancyhdr}
\usepackage{xcolor}

\hypersetup{
  colorlinks=true,
  linkcolor=blue!50!black,
  citecolor=blue!50!black,
  urlcolor=blue!50!black
}

\captionsetup{font=small}

\newcommand{\TitreRapport}{Classification de nuages de points 3D avec PointNet}
\newcommand{\SousTitreRapport}{Master PAR --- Lab PointNet}
\newcommand{\AuteurRapport}{NOCHI MAGOUO}
\newcommand{\Encadrants}{Jean-Emmanuel DESCHAUD \\ Paul CHECCHIN}
\newcommand{\UE}{3D Deep Learning}
\newcommand{\Annee}{2025--2026}

\newcommand{\fairepagedegarde}{%
\begin{titlepage}
\centering
\vspace*{0.8cm}

\includegraphics[width=0.22\textwidth]{report_assets/ensias.png}\par
\vspace{0.7cm}
{\Large Ecole Universitaire de Physique et d'Ingenierie\par}
\vspace{0.25cm}
{\large \UE\par}
\vspace{0.2cm}
{\large \Annee\par}

\vfill
{\Large \textbf{Rapport de projet}}\par
\vspace{0.4cm}
{\huge \textbf{\TitreRapport}}\par
\vspace{0.3cm}
{\Large \SousTitreRapport}\par
\vspace{0.7cm}
\rule{0.88\textwidth}{0.5pt}\par
\vspace{0.9cm}

\begin{tabular}{>{\raggedright\arraybackslash}p{0.44\textwidth} >{\raggedright\arraybackslash}p{0.44\textwidth}}
\textbf{Realise par} & \textbf{Encadrants} \\
\AuteurRapport & \Encadrants \\
\end{tabular}

\vfill
{\large \today\par}
\end{titlepage}
}

\pagestyle{fancy}
\fancyhf{}
\lhead{PointNet Lab --- M2 PAR}
\rhead{\nouppercase{\leftmark}}
\cfoot{\thepage}
\setlength{\headheight}{14pt}

\begin{document}

\hypersetup{pageanchor=false}
\fairepagedegarde

\pagenumbering{roman}
\tableofcontents
\listoffigures
\listoftables
\clearpage
\pagenumbering{arabic}
\hypersetup{pageanchor=true}

\chapter*{Resume}
\addcontentsline{toc}{chapter}{Resume}
Ce rapport presente l'implementation en PyTorch des exercices PointMLP, PointNetBasic et PointNetFull (T-Net 3\(\times\)3), puis l'etude de plusieurs strategies de data augmentation sur ModelNet10.
La campagne de reference du notebook final atteint \textbf{90.09\%} d'accuracy test avec la configuration \emph{baseline + occlusion}.
Les campagnes automatisees confirment que l'occlusion reste l'une des augmentations les plus robustes selon les metriques de validation/evaluation disponibles.

\chapter{Introduction}

\section{Contexte du module}
Dans ce module, je travaille sur le traitement de nuages de points 3D.
Un point cloud est un ensemble de points \((x,y,z)\) non ordonnes: je ne peux donc pas supposer un voisinage fixe comme dans une image 2D.
Le probleme principal est de construire un modele qui respecte cette structure en ensemble.

\section{Lien avec l'article PointNet}
Le projet est inspire de l'article \textit{PointNet: Deep Learning on Point Sets for 3D Classification and Segmentation}~\cite{qi2017pointnet}.
Avant d'implementer mes reseaux, j'ai essaye de comprendre sa logique:
\begin{itemize}
  \item traiter directement les points, sans voxelisation lourde;
  \item apprendre des descripteurs locaux par point via des MLP partages;
  \item utiliser une fonction symetrique (max-pooling) pour obtenir une invariance a l'ordre des points;
  \item ajouter un T-Net pour stabiliser l'alignement geometrique en entree.
\end{itemize}

\section{Objectif de mon rapport}
Je veux montrer que je comprends ce que je fais, et pas seulement que j'obtiens un score.
Concretement, je documente:
\begin{itemize}
  \item mon evolution de PointMLP vers PointNet;
  \item ma methode experimentale (runs, controles, comparaison);
  \item pourquoi j'ai retenu \texttt{baseline + occlusion} en fin de projet.
\end{itemize}

\chapter{Methodes et architectures}

\section{Idee generale que je retiens de PointNet}
La forme que je retiens de l'article est:
\begin{align}
f(\{x_1, \dots, x_n\}) \approx \gamma\left(\max_i\, h(x_i)\right)
\end{align}
avec \(h\) un encodeur partage par point et \(\max\) une aggregation symetrique.
Cette ecriture m'a servi de fil conducteur pour comprendre pourquoi PointNet fonctionne mieux qu'un MLP classique sur des donnees non ordonnees.

\section{Reseau 1: PointMLP (Exercice 1)}
J'ai d'abord implemente le modele demande dans l'enonce: flatten \(1024\times 3\) puis couches fully-connected.

\begin{figure}[H]
\centering
\begin{tikzpicture}[node distance=0.75cm]
    \node[input] (in) {Input Point Cloud \\ $(N \times 3)$};
    \node[layer, below=of in] (flatten) {Flatten \\ $(3072)$};
    \node[layer, below=of flatten] (fc1) {FC + BN + ReLU \\ $(3072 \rightarrow 512)$};
    \node[layer, below=of fc1] (fc2) {FC + BN + ReLU \\ $(512 \rightarrow 256)$};
    \node[layer, below=of fc2] (drop) {Dropout \\ $(p=0.3)$};
    \node[layer, below=of drop] (fc3) {FC \\ $(256 \rightarrow K)$};
    \node[layer, below=of fc3] (logsoft) {LogSoftmax};
    \node[input, fill=red!10, below=of logsoft] (out) {Output Classes \\ $(K)$};

    \draw[arrow] (in) -- (flatten);
    \draw[arrow] (flatten) -- (fc1);
    \draw[arrow] (fc1) -- (fc2);
    \draw[arrow] (fc2) -- (drop);
    \draw[arrow] (drop) -- (fc3);
    \draw[arrow] (fc3) -- (logsoft);
    \draw[arrow] (logsoft) -- (out);
\end{tikzpicture}
\caption{Architecture du PointMLP (conforme au notebook: 3072\(\rightarrow\)512\(\rightarrow\)256\(\rightarrow\)K).}
\label{fig:arch-mlp}
\end{figure}

\noindent\textbf{Analyse.} Avec ce modele, je perds la structure "ensemble de points" au moment du flatten.
Le modele voit un grand vecteur ordonne, alors que l'ordre des points est arbitraire.
C'est la limite qui m'a pousse a passer a PointNetBasic.

\section{Reseau 2: PointNetBasic (Exercice 2.1)}
Dans PointNetBasic, je remplace le flatten par des MLP partages (\texttt{Conv1d} de noyau 1) puis un max-pooling global.

\begin{figure}[H]
\centering
\resizebox{0.98\linewidth}{!}{%
\begin{tikzpicture}[node distance=0.6cm and 1cm]
    \node[input] (in) {Input Cloud \\ $(B, 3, N)$};

    \node[layer, right=of in] (conv1) {Conv1D(64) + BN + ReLU};
    \node[layer, right=of conv1] (conv2) {Conv1D(64) + BN + ReLU};
    \node[layer, right=of conv2] (conv3) {Conv1D(64) + BN + ReLU};

    \node[layer, below=of conv3] (conv4) {Conv1D(128) + BN + ReLU};
    \node[layer, left=of conv4] (conv5) {Conv1D(1024) + BN + ReLU};

    \node[pool, left=of conv5] (pool) {Global Max Pooling \\ $(B, 1024, 1)$};
    \node[layer, below=of pool] (fc1) {FC(512) + BN + ReLU};
    \node[layer, right=of fc1] (fc2) {FC(256) + BN + ReLU};
    \node[layer, right=of fc2] (drop) {Dropout};
    \node[layer, right=of drop] (fc3) {FC(K) + LogSoftmax};

    \node[input, fill=red!10, below=of fc3] (out) {Output $(B, K)$};

    \draw[arrow] (in) -- (conv1);
    \draw[arrow] (conv1) -- (conv2);
    \draw[arrow] (conv2) -- (conv3);
    \draw[arrow] (conv3) -- (conv4);
    \draw[arrow] (conv4) -- (conv5);
    \draw[arrow] (conv5) -- (pool);
    \draw[arrow] (pool) -- (fc1);
    \draw[arrow] (fc1) -- (fc2);
    \draw[arrow] (fc2) -- (drop);
    \draw[arrow] (drop) -- (fc3);
    \draw[arrow] (fc3) -- (out);
\end{tikzpicture}
}
\caption{Architecture du PointNetBasic (sans T-Net), basee sur les blocs \texttt{Conv1d} du notebook.}
\label{fig:arch-basic}
\end{figure}

\noindent\textbf{Analyse.} Ce schema corrige le probleme principal de PointMLP:
le max-pooling rend la sortie invariante a la permutation des points.
Pour le traitement de point clouds, c'est une difference fondamentale.

\section{Reseau 3: PointNetFull avec T-Net 3x3 (Exercice 2.2)}
La version Full ajoute un T-Net d'alignement 3x3 en entree.
J'utilise aussi une regularisation orthogonale:
\begin{align}
\mathcal{L} = \mathcal{L}_{CE} + \lambda\lVert I-AA^\top\rVert_F^2
\end{align}

\begin{figure}[H]
\centering
\resizebox{0.98\linewidth}{!}{%
\begin{tikzpicture}[node distance=0.8cm and 1.5cm]
    \node[input] (in) {Input Cloud \\ $(B, N, 3)$};

    \node[tnet, below right=of in, yshift=1cm] (tnet) {Tnet ($3 \times 3$)};
    \node[op, right=of in] (mul) {$\otimes$};

    \node[layer, right=of mul] (features) {PointNetBasic \\ Feature Extraction \\ (Conv1D jusqu'a MaxPool)};
    \node[layer, right=of features] (mlp) {MLP Head \\ (FC 512 \(\rightarrow\) 256 \(\rightarrow\) K)};
    \node[input, fill=red!10, right=of mlp] (out) {Scores \\ $(B, K)$};

    \draw[arrow] (in) -- (mul);
    \draw[arrow] (in) |- (tnet);
    \draw[arrow] (tnet) -| node[pos=0.25, right] {Matrice de transformation ($3 \times 3$)} (mul);
    \draw[arrow] (mul) -- node[above] {$(B, 3, N)$} (features);
    \draw[arrow] (features) -- node[above] {Global Feature $(B, 1024)$} (mlp);
    \draw[arrow] (mlp) -- (out);
\end{tikzpicture}
}
\caption{Architecture du PointNetFull avec module \texttt{Tnet} d'alignement en entree.}
\label{fig:arch-full}
\end{figure}

\noindent\textbf{Analyse.} Le T-Net apporte une adaptation geometrique supplementaire.
Dans ma lecture, il agit comme une normalisation apprenable: le reseau apprend a "mettre en forme" le nuage avant extraction des features globales.

\section{Comparaison structurelle des trois reseaux}
\begin{table}[H]
\centering
\caption{Comparaison de mes trois implementations.}
\label{tab:archi-compare}
\begin{tabular}{lccc}
\toprule
Modele & Invariance ordre & Alignement geometrique & Params trainables \\
\midrule
PointMLP & Non & Non & 1,708,810 \\
PointNetBasic & Oui (max-pooling) & Non & 811,914 \\
PointNetFull & Oui (max-pooling) & Oui (T-Net 3x3) & 1,614,995 \\
\bottomrule
\end{tabular}
\end{table}

\noindent\textbf{Analyse.} J'insiste sur un point: l'amelioration ne vient pas simplement d'avoir plus de parametres.
Elle vient surtout du fait que PointNetBasic et PointNetFull respectent mieux la nature d'un point cloud (ensemble non ordonne, variations geometriques).

\chapter{Protocole experimental}

\section{Jeux de donnees et splits}
Le notebook final utilise:
\begin{itemize}
  \item train: \textbf{3193} exemples,
  \item validation: \textbf{798} exemples,
  \item test: \textbf{908} exemples.
\end{itemize}
Le split est derive du train officiel avec \texttt{VAL\_RATIO = 0.2} et \texttt{seed = 42}, puis le test officiel est conserve pour l'evaluation finale.

\section{Configuration du notebook final}
Configuration extraite de \texttt{Code/pointnet\_notebook.ipynb}:
\begin{itemize}
  \item modele: \texttt{PointNetFull},
  \item augmentations train: \texttt{[normalize, rot\_z, noise, shuffle, occlusion]},
  \item augmentation test: \texttt{[normalize]},
  \item \texttt{MAX\_EPOCHS = 150},
  \item early stopping active (monitor \texttt{val\_acc}, patience \texttt{50}, min\_delta \texttt{0.0}),
  \item scheduler de taux d'apprentissage actif.
\end{itemize}

\section{Campagnes automatisees complementaires}
En parallele du notebook, des campagnes de benchmark ont ete lancees avec le script robuste.
Table~\ref{tab:campaigns} resume les runs principaux exploites dans ce rapport.

\begin{table}[H]
\centering
\caption{Campagnes experimentales utilisees.}
\label{tab:campaigns}
\begin{tabular}{p{0.22\linewidth}p{0.15\linewidth}p{0.57\linewidth}}
\toprule
Campagne & Date & Usage dans le rapport \\
\midrule
C1 (20260208\_145104) & 08/02/2026 & benchmark de 10 augmentations (seed 42, 25 epochs) \\
C2 (20260211\_081948) & 11/02/2026 & benchmark robuste (4 experiences, seed 1, max 50 epochs, early stop) \\
C3 (20260211\_143359) & 11/02/2026 & run cible \texttt{plus\_occlusion, seed\_0} (max 150 epochs, monitor val\_loss) \\
\bottomrule
\end{tabular}
\end{table}

\section{Metriques reportees}
Selon les sorties disponibles, on suit principalement:
\begin{itemize}
  \item \texttt{best\_val\_acc} et \texttt{best\_val\_loss},
  \item \texttt{best\_eval\_acc} quand presente,
  \item \texttt{final\_test\_acc} pour le notebook final,
  \item matrices de confusion et courbes d'apprentissage.
\end{itemize}
Cette distinction est importante pour eviter de melanger validation interne et evaluation finale.

\chapter{Resultats}

\section{Campagne augmentation seed 42 (25 epochs)}
Le benchmark du 8 fevrier 2026 compare 10 configurations autour de PointNetFull.
Les meilleures configurations en \texttt{best\_val\_acc} sont presentees dans le tableau~\ref{tab:aug-20260208}.

\begin{table}[H]
\centering
\caption{Top configurations --- campagne \texttt{20260208\_145104}.}
\label{tab:aug-20260208}
\begin{tabular}{lccc}
\toprule
Configuration & best\_val\_acc (\%) & best\_epoch & final\_val\_acc (\%) \\
\midrule
plus\_jitter & 91.08 & 23 & 86.67 \\
plus\_mirror & 90.75 & 21 & 87.56 \\
plus\_occlusion & 90.42 & 23 & 89.76 \\
baseline & 90.31 & 23 & 88.77 \\
combo\_full & 90.09 & 22 & 90.09 \\
\bottomrule
\end{tabular}
\end{table}

\begin{figure}[H]
  \centering
  \includegraphics[width=0.82\linewidth]{../Code/reports/pointnetfull_aug_suite_20260208_145104/figures/global_best_val_accuracy_bar.png}
  \caption{Comparaison globale des \texttt{best\_val\_acc} (campagne du 08/02/2026).}
  \label{fig:global-bar-20260208}
\end{figure}

\section{Benchmark robuste (seed 1, max 50 epochs)}
Le benchmark du 11 fevrier 2026 confirme que \texttt{plus\_occlusion} reste en tete parmi les 4 configurations testees (Table~\ref{tab:robust-20260211}).

\begin{table}[H]
\centering
\caption{Synthese --- benchmark robuste \texttt{20260211\_081948}.}
\label{tab:robust-20260211}
\begin{tabular}{lccc}
\toprule
Configuration & best\_val\_acc (\%) & best\_epoch & final\_val\_acc (\%) \\
\midrule
plus\_occlusion & 88.33 & 21 & 87.33 \\
combo\_full & 88.22 & 18 & 88.00 \\
baseline & 86.45 & 10 & 81.61 \\
plus\_jitter & 86.12 & 7 & 83.81 \\
\bottomrule
\end{tabular}
\end{table}

\section{Run de reference du notebook final}
Le run final du notebook (seed 42) donne:
\begin{itemize}
  \item \texttt{MAX\_VAL\_ACC = 94.36\%} (sur 150 epochs),
  \item \texttt{Final TEST accuracy = 90.09\%},
  \item \texttt{Final TEST loss = 0.4077}.
\end{itemize}
Ces valeurs sont celles affichees directement dans la sortie du notebook final.

\section{Run cible occlusion (seed 0)}
Dans la campagne \texttt{20260211\_143359}, le run \texttt{plus\_occlusion/seed\_0} atteint:
\begin{itemize}
  \item \texttt{best\_val\_acc = 90.09\%} (epoch 23),
  \item \texttt{best\_eval\_acc = 89.43\%},
  \item \texttt{final\_val\_acc = 88.55\%} apres early stopping.
\end{itemize}

\begin{figure}[H]
  \centering
  \includegraphics[width=0.82\linewidth]{../Code/reports/pointnetfull_aug_suite_20260211_143359/runs/plus_occlusion/seed_0/curves.png}
  \caption{Courbes d'entrainement du run \texttt{plus\_occlusion, seed\_0}.}
  \label{fig:occlusion-seed0-curves}
\end{figure}

\begin{figure}[H]
  \centering
  \begin{subfigure}{0.48\linewidth}
    \centering
    \includegraphics[width=\linewidth]{../Code/reports/pointnetfull_aug_suite_20260208_145104/figures/baseline_confusion_matrix.png}
    \caption{Baseline}
  \end{subfigure}
  \hfill
  \begin{subfigure}{0.48\linewidth}
    \centering
    \includegraphics[width=\linewidth]{../Code/reports/pointnetfull_aug_suite_20260211_143359/runs/plus_occlusion/seed_0/confusion_matrix.png}
    \caption{Plus occlusion (seed 0)}
  \end{subfigure}
  \caption{Comparaison qualitative des matrices de confusion.}
  \label{fig:cm-compare}
\end{figure}

\chapter{Discussion}

\section{Lectures principales}
Trois constats ressortent des campagnes disponibles:
\begin{itemize}
  \item l'\textbf{occlusion} est systematiquement competitive et souvent premiere dans les protocoles les plus contraints;
  \item les combinaisons lourdes d'augmentations ne garantissent pas de gain stable par rapport a une politique plus simple;
  \item la configuration du monitor d'early stopping (\texttt{val\_acc} vs \texttt{val\_loss}) influence notablement l'epoch retenue et les metriques finales.
\end{itemize}

\section{Interpretation}
L'occlusion aleatoire est pertinente pour ModelNet10 car elle force le reseau a reconnaitre des formes partielles.
Cela ameliore la robustesse a la disparition locale de points, ce qui est coherent avec des scenes reelles ou les acquisitions sont incompletes.
Le gain en test du notebook final (\(90.09\%\)) valide ce choix pour la version courante du projet.

\section{Limites actuelles}
\begin{itemize}
  \item les benchmarks robustes ne couvrent pas encore un vrai multi-seeds complet (plusieurs tables ont \texttt{num\_runs = 1});
  \item la comparaison Exercice 1 / Exercice 2 reste surtout architecturale dans cette version du rapport, faute de tableaux consolides homogènes sur la meme grille d'essais;
  \item l'evaluation est concentree sur ModelNet10, sans extension ModelNet40.
\end{itemize}

\section{Etat en cours}
Au moment de redaction, un fine-tuning \texttt{baseline + occlusion} en \texttt{seed 45} est en cours.
Il est prevu d'ajouter ses metriques finales (accuracy, loss, matrice de confusion) dans le tableau de comparaison principal des que le run est termine.

\chapter{Conclusion}
Ce projet a permis d'implémenter PointMLP et PointNet en PyTorch, puis de construire un protocole expérimental propre
(train/val/test séparés, early stopping sur validation, test final unique).

Le résultat final atteint environ \textbf{90\%} d'accuracy test sur ModelNet10 avec une configuration
\textbf{baseline + occlusion}, ce qui constitue un compromis performant et lisible pour ce TP.


\clearpage
\addcontentsline{toc}{chapter}{Bibliographie}
\bibliographystyle{plain}
\bibliography{references}

\end{document}
